\chapter{Repeated collections decisions and policy value}\label{ch:collections}
\section{Define assistance before optimizing it}
Collections is a natural longitudinal application because an action changes
the state that determines later action. A payment reminder, temporary
forbearance and a contractual restructuring are different interventions.
They affect cash timing, costs, eligibility and customer outcomes differently.
The initial product should evaluate a small set of operationally defined
assistance strategies, with ordinary servicing as the comparator, before
attempting unrestricted policy optimization.

The economic owner matters. A servicer, a loan owner and a purchaser of
distressed debt receive different cash flows. The study preserves the
original eligible cohort, includes assistance delivery costs and defines
the remaining claim at the horizon. It reports customer outcomes separately,
such as persistent hardship or repeated delinquency, rather than assuming
that bank contribution fully represents them.

The following complete synthetic study isolates the observed-history case.
The controlled hidden-state study in Chapter~\ref{ch:latentestimation}
then shows what changes when relevant information is hidden from the analyst.
Neither study represents a named bank's customers or treatment effectiveness.

\section{A three-state decision process}
The observed pre-action state is $S_t\in\{0,1,2\}$, representing
current, moderate stress and severe stress. Initial proportions are
$(0.6,0.3,0.1)$. Extra assistance is feasible only in states one and two.
The historical action-one probabilities are $(0,0.35,0.60)$.
These are known randomized probabilities conditional on the observed state;
there is no hidden action confounding in this experiment.

The transition matrices are
\[
 P_0=\begin{pmatrix}.88&.10&.02\\.30&.50&.20\\.08&.22&.70\end{pmatrix},
 \qquad
 P_1=\begin{pmatrix}.88&.10&.02\\.50&.40&.10\\.18&.42&.40\end{pmatrix}.
\]
The first row for action one is not used because assistance is infeasible
there. The process has eight decisions and monthly discount $\gamma=0.99$.
Contribution associated with the next state is $(20,-10,-160)$ dollars,
less USD\,12 for each assistance action. These are signed monthly
contribution increments, not a repeated write-off of the same principal.
A terminal adjustment of USD\,$-100$ in severe stress represents the
remaining account beyond the eight-month cash window.

This law lets assistance improve future states while costing money today.
It also creates practical support difficulties: the probability of following
a deterministic strategy over all eight dates can be small even though
each feasible action has positive probability at every state. The known
finite Markov law gives an exact policy-value reference without requiring
Monte Carlo approximation of truth.

\section{Policies, training information and exact evaluation}
Three fixed policies are declared before fitting: no extra assistance,
assistance only in severe stress, and assistance in all eligible states.
They use observable information and respect the same feasible action set.
The study draws 3,000 training trajectories and 8,000 independent evaluation
trajectories. Training data estimate transition rows by observed counts
with a one-half pseudocount in each next-state category.

For each target policy $\pi$, backward recursion under the fitted
transition law gives
\[
 \widehat Q_t(s,a)=\sum_{s'}\widehat P_a(s,s')
 \{r(s')+\gamma\widehat V_{t+1}(s')\}-12a,
 \quad \widehat V_t(s)=\sum_a\pi(a\mid s)\widehat Q_t(s,a).
\]
The terminal condition is the declared remaining account. Repeating the
recursion with the generating transition law gives the exact reference
value, available only to the evaluator. The fitted policy scores never
consume generated latent variables or future information unavailable at
the relevant decision.

The direct method averages $\widehat V_0(S_0)$ over evaluation initial
states. Stepwise importance sampling reweights each dated reward by the
full product of target-to-behavior action ratios through that date.
Sequential DR adds cumulative-ratio Bellman residuals to the direct
prediction. All three target the same discounted contribution and
terminal account. This alignment is essential for a fair comparison.

\begin{table}[htbp]\centering\small
\begin{tabular}{@{}lrrrrr@{}}
\toprule
Policy & True USD & Direct & IPS & DR & Final ESS\\
\midrule
No extra assistance & -139.58 & -123.44 & -135.08 & -137.56 & 146\\
Severe state only & -58.26 & -49.57 & -59.95 & -59.55 & 1267\\
All eligible & -2.34 & 0.48 & 2.93 & 1.70 & 489\\
\bottomrule
\end{tabular}

\caption{Eight-decision fixed-policy values in dollars per original eligible
customer. Final ESS describes concentration of final trajectory weights;
it is not the independent sample size of every dated score.}
\label{tab:collections-values}
\end{table}

The direct estimates depend entirely on the fitted transition law.
Importance sampling uses known logging and remains highly variable on
rare complete paths. DR corrects fitted continuation errors but cannot
remove all realized transition and reward uncertainty. Its empirical
value need not lie between the other two estimates.

\section{A plausible shortcut estimates the wrong law}
Suppose the analyst uses only the current-action ratio at each date.
This changes the distribution of the action conditional on the historical
incoming state, but leaves that incoming state distributed under historical
practice. Earlier assistance has already changed who reaches severe
stress. The shortcut therefore does not evaluate a sustained target policy.

For the all-eligible strategy in this run, that current-ratio calculation
gives approximately USD\,$-36.79$, compared with exact value USD\,$-2.34$.
For no extra assistance it gives about USD\,$-82.67$, compared with exact
value USD\,$-139.58$. The error is not fixed by a larger evaluation sample:
it is a mismatch between the reweighted distribution and the requested
intervention. A marginalized state-ratio estimator would need the target-
to-behavior state occupancy ratio as well as the action ratio.

The software consequence is a required weight-scope declaration.
``Propensity-weighted'' is too vague. A result should distinguish
current-action, full-history, marginal occupancy and censoring weights,
with their target and normalization. A consumer can then reject a score
that lacks the history correction required by its policy contrast.

\section{Evaluate contrasts with their covariance}
The practical question is the incremental contribution relative to the
incumbent. Compute both policy scores on each evaluation trajectory and
form their difference. The standard error of the mean uses the variance
of this paired difference, preserving covariance from shared customers.
There are two prespecified comparisons against no extra assistance;
Bonferroni-adjusted normal intervals give a conservative simultaneous
approximation for that family.

\begin{table}[htbp]\centering\small
\begin{tabular}{@{}lrrr@{}}
\toprule
Policy versus no extra & True USD & DR difference & Simultaneous interval\\
\midrule
Severe state only & 81.32 & 78.01 & [44.58, 111.44]\\
All eligible & 137.24 & 139.26 & [103.83, 174.69]\\
\bottomrule
\end{tabular}

\caption{Prespecified contrasts using paired DR scores. These intervals
condition on the training fits and on the synthetic complete-trajectory law.}
\label{tab:collections-contrasts}
\end{table}

The account and intervention are deliberately simple enough to reconstruct.
For a bank implementation, inference must also handle estimated observation
probabilities, household or regional clustering, delayed losses and policy
selection. The positive synthetic comparisons cannot be treated as measured
collections benefit for the investment case.

\section{Double robustness and interval reliability differ}
The repeated experiment uses 160 independent datasets, each with 600
training and 2,000 evaluation trajectories. It evaluates the all-eligible
policy under four nuisance specifications: fitted transition values and
true logging; zero interior action values and true logging; exact action
values and a wrong logging model assigning probability $0.5$ in eligible
states; and both the zero action values and wrong logging.

The terminal account is correct in all four cases. The first two illustrate
the protection of known assignment; the third isolates the correct-value
route; the fourth has no double-robustness protection. The generating law
and declared policies remain fixed across repetitions, while the independent
datasets vary. No favorable subset of repetitions is selected.

\begin{table}[htbp]\centering\small
\begin{tabular}{@{}lrrr@{}}
\toprule
Specification & Bias (USD) & RMSE & 95\% coverage\\
\midrule
DR & 0.15 & 17.33 & 0.938\\
Wrong Q; true logging & 0.71 & 18.28 & 0.869\\
True Q; wrong logging & 0.55 & 10.72 & 0.912\\
Both wrong & 1.41 & 12.67 & 0.856\\
\bottomrule
\end{tabular}

\caption{Repeated-sample error and nominal normal-interval coverage for
all-eligible policy value. Monte Carlo standard errors for coverage range
from about 0.019 to 0.028 and are retained in the result file.}
\label{tab:collections-robustness}
\end{table}

The zero-value/true-logging case has mean protection but coverage well below
95\% in this finite experiment. Highly variable path weights make a
normal approximation unreliable at this sample size. Even the exact-value
case does not guarantee nominal finite-sample coverage. The wrong-both
case happens to have modest average bias under this particular law; that
is not evidence of a new robustness property. A misspecified estimator can
perform tolerably under one mechanism and fail under another.

The product recommendation is therefore to report weight concentration,
extreme score contributions, horizon sensitivity and repeated-sample
calibration on relevant synthetic mechanisms before selecting an interval
method. Weight truncation can improve stability while introducing bias;
any chosen threshold is part of the estimator. A shorter feasible policy
comparison or a prospective experiment may be more informative than a
nominally sophisticated long-horizon estimate.

\section{What partial observation adds to this study}
Observed stress states in the present experiment are sufficient for the
transition and action laws. Real collections staff may observe private
conversation details, perceived ability to pay or impending income events.
If those signals also affect action assignment and outcomes, the
recorded-history exchangeability used here fails.

Chapter~\ref{ch:latentestimation} constructs that alternative explicitly.
Signals imperfectly measure a hidden state, and historical assistance
depends on that state. A joint action--observation likelihood improves
recovery under the stated parametric environment but retains substantial
policy uncertainty. Proximal and instrumental-variable OPE offer other
routes under different measurements and restrictions
\citep{tennenholtz2020,bennettprox2023,xu2023}. A bank cannot choose
among those routes solely by whichever model gives the largest estimated
recovery gain.

The first bank pilot should therefore investigate the action log alongside
the outcome data. It must document which staff information was available,
which variables were recorded before decisions, how discretion operated,
and whether an existing randomized or quasi-experimental variation can
anchor transition effects. That work determines whether the module returns
an identified policy contrast, a conditional scenario or a bounded effect.

\section{A deployable study application, before an automated policy}
The standalone collections workbench would display complete eligible
trajectories, action and observation clocks, fitted continuation values,
support diagnostics, fixed-policy comparisons and customer-outcome
constraints. It would retain the original cohort after default, sale or
closure and explain how each event enters the economic account. A model
developer can investigate latent dynamics without promoting a policy.

A downstream application can consume the evidence package to compare
approved strategies. Promotion to operational decision making additionally
requires current permissions, execution capability, customer protection,
prospective or otherwise adequate value evidence and ongoing monitoring.
Those are implementation obligations of the proposed bank product, not
claims established by the present synthetic studies.
